% ===========================================================================
\section{The Correlated-Factor Case}
\label{app:uncertainty}
% ===========================================================================

This table is a retained reference-year diagnostic of two declared fuel-factor dependence settings. It uses the stated legacy seed and parameter implementation; it is distinct from the current common annual conditional ensemble, which shares named physical inputs and pairs years within each current dependence case. The diagnostic does not define a mathematical uncertainty bound.

\begin{table}[htbp]
\centering
\footnotesize
\begin{tabular}{@{}lrr@{}}
\toprule
& \textbf{Independent} & \textbf{Correlated} \\
\midrule
Sector mean (Mt \COe{}) & 41.926 & 41.922 \\
2.5th percentile        & 36.869 & 36.891 \\
97.5th percentile       & 47.638 & 47.642 \\
\textbf{Sector $\pm$\%} & \textbf{12.9} & \textbf{12.9} \\
\midrule
Road transport $\pm$\%     & 13.2 & 13.2 \\
Mixed off-road $\pm$\%     & 29.0 & 28.9 \\
\bottomrule
\end{tabular}
\caption{Approach~2, \num{10000} realisations at seed \num{42}, from
\texttt{outputs/national\_uncertainty\_montecarlo.csv}, retained as the legacy reference diagnostic. The central estimate does not
move materially---correlation does not bias a sum---and under the common-random-numbers
scheme the band does not move either: both treatments print the same width ($\pm\SI{12.9}{\percent}$).}
\label{tab:mc}
\end{table}

\subsection{What the two scenarios assume}

Five sector modules in the national inventory burn fossil fuel, and they draw their
carbon-content factors from the same IPCC table entries with the same published uncertainty
ranges. The \emph{independent} scenario draws each sector's factor separately. The
\emph{correlated} scenario draws one realisation of each fuel's factor and applies it
everywhere at once.

The correlated case is physically the more defensible of the two extremes. There is one carbon
content of the diesel Pakistan burns, and if the IPCC default is wrong about it, it is wrong
in the same direction for road transport, railways, off-road machinery and captive generation
simultaneously. Treating the draws as independent asserts that the true value could be at the
bottom of its range for railways and the top for road transport at the same moment, which it
cannot.

\subsection{Why this sector no longer moves}

The retained reference diagnostic rounds to $\pm\SI{13.2}{\percent}$ under either setting---the
common-random-numbers rebuild made the correlation scenario a no-op here. Mixed off-road
machinery likewise, at $\pm\SI{28.9}{\percent}$ against $\pm\SI{28.9}{\percent}$.

The difference between the two is instructive. Road transport's uncertainty is dominated by
the shared fuel-carbon factor, because everything else about it is comparatively tight: the
activity datum is a metered national supply, and \COtwo{} is \SI{97.5}{\percent} of the
emissions. When the factor stops being able to average against other sectors' draws, there is
very little else left to average, so almost the whole band is exposed. Mixed off-road machinery
carries a large independent uncertainty of its own---on the energy share that
carves it out, and on the \NtwoO{} factor---which continues to average regardless, so the
correlation adds proportionally little. Note also why the sector band sits below both
category bands in either scenario: the partition uncertainty between road and off-road is
negatively coupled and cancels on recombination, while the shared control and
carbon-factor terms are positively correlated and survive.

Anchoring is what produces this shape. An inventory that built road transport bottom-up would
carry substantial independent uncertainty in its fleet model, and that uncertainty would
partly mask the factor correlation. Anchoring removes the fleet uncertainty from the total and
leaves the factor exposed, which is a more honest representation of what is actually not
known.

\subsection{How to read the sector's band}

The earlier sector comparisons below retain their own declared parameter settings and are not the common annual ensemble. The current release samples shared fuel-carbon identities consistently across GHG sources and pairs years within each named dependence case. Perfect persistence and annually varying activity errors are sensitivity assumptions rather than mathematical bounds. Parameter percentiles remain conditional on the road/off-road claimant assignment, transferred factors and admitted sources. Class, sulfur and technology endpoints are structural alternatives and receive no probability.

Beside the correlated-factor case, the sector publishes its own structural-scenario
table (\texttt{outputs/transport\_uncertainty\_scenarios\_2024.csv}): the boundary cases
of the main paper's boundary table (the national-minus-bulk construction, the road-share
envelope, the retail non-road cases, the calendar-year construction), the provenance-audit
bridges beside it, the CNG heating-value case, the off-road \NtwoO{}
table choice and share cases, the factor-mix invariance row, the CNG exclusion bound, the
fleet-basis energy band, the off-road reporting recode, the class-composition cases at the
ends of the declared MJ/km bands and the contrast cases that move light and heavy classes
against one another, and two declaration rows---the class split and the air-factor
provenance---that are scope decisions rather than samplable parameters. Each carries a conditional parametric interval, an explicit
declaration, or a proved invariance, never a structural probability. The deduction and
cross-check inputs each carry their own register row: the generator carve-out (survey
expenditure over a declared administered price), the off-road energy fraction,
and the rail-figure discrepancies---\SI{3.8}{\percent} year-matched,
\SI{11.2}{\percent} cross-year---that anchor the controls' declared \SI{10}{\percent}
(\texttt{outputs/transport\_uncertainty\_inputs\_2024.csv}).
